\documentclass[11pt]{article}
\usepackage{problem_set_style}
\usepackage{kotex}
\usepackage{booktabs}
\usepackage{array}
\begin{document}
\heading{}{샘플}{}{Spring 2027}{프로젝트 1 · 동료평가 샘플 (8주차 제출)}

\noindent{\small 아래는 \textbf{가상의 팀}이 쓴 샘플입니다. 수치는 형식을 보이기 위한 예시이며 실제 실험 결과가 아닙니다. 형식과 분량의 기준으로만 참고하세요.}

\noindent\textbf{평가 대상}: 다른 팀의 ``당뇨 재입원 예측'' 보고서

\begin{center}\small
\begin{tabular}{@{}p{3.4cm}p{9cm}@{}}\toprule
항목 & 평가 \\\midrule
주장의 명확성 & 첫 문단에서 입원 한 건의 30일 내 재입원을 예측한다가 분명하다. 라벨과 지표도 명확하다. \\
비교 대상 & 항상 재입원 없음을 내놓는 베이스라인과 로지스틱회귀, GBDT를 같은 분할로 비교했다. 좋다. \\
확장성 & 한 병원의 데이터만 썼다. 다른 병원이나 다른 연도에서도 통하는지 확인이 없다. \\
재현성 & 시드와 하이퍼파라미터는 있으나, 한 환자의 여러 입원을 어떻게 나눴는지와 코드 실행 방법이 적혀 있지 않아 다시 해 볼 수 없다. \\
한계 & 실패한 시도와 클래스 불균형의 영향을 솔직하게 적었다. 좋다. \\
\bottomrule\end{tabular}\end{center}

\noindent\textbf{반박 가능한 질문.} 같은 환자의 입원이 train과 test에 함께 들어갔다면 test 숫자는 부풀려진 것이 아닐까요? 환자 단위로 나눠서 다시 측정한 결과가 있나요?

\end{document}
